% ════════════
% WP5_Working_Paper.tex
% Relational Engineering Specifications: The Operational Architecture of the Warm Offer, Social Network Transition Protocol, and Deployment Counterpoint Framework
% Material Dignity Infrastructure Working Paper 5
% ════════════

\documentclass[12pt, letterpaper]{article}

% ── PACKAGES ──
\usepackage[left=0.7in, right=0.7in, top=1in, bottom=1in, headheight=14pt]{geometry}
\usepackage[expansion=false]{microtype}
\usepackage{textcomp}
\usepackage{setspace}
\usepackage{booktabs}
\usepackage{tabularx}
\usepackage[titles]{tocloft}
\usepackage{tabularray}
\usepackage{longtable}
\usepackage{array}
\usepackage{multirow}
\usepackage{hyperref}
% natbib not used — manual thebibliography
\usepackage{amsmath}
\usepackage{amssymb}
\usepackage{parskip}
\usepackage{titlesec}
\usepackage{fancyhdr}
\usepackage[table]{xcolor}
\usepackage{caption}
\usepackage{footnote}
\usepackage{enumitem}
\usepackage{tikz}
\usepackage{needspace}
\usepackage{etoolbox}
\usepackage{float}

% ── COLOR SYSTEM ───
\definecolor{auditblue}{HTML}{003366}
\definecolor{rowgray}{HTML}{F5F5F5}

% ── GLOBAL SPACING ───
\setstretch{1.4}
\setlength{\parindent}{0pt}
\setlength{\parskip}{8pt}
\hyphenpenalty=1000

% ── SECTION FORMATTING. ZERO BOLDING IN TEXT. ───
\titleformat{\section}
  {\normalfont\large\color{auditblue}}{\thesection.}{0.5em}{}
\titleformat{\subsection}
  {\normalfont\normalsize\color{auditblue}}{\thesubsection.}{0.5em}{}
\titlespacing*{\section}{0pt}{18pt}{6pt}
\titlespacing*{\subsection}{0pt}{12pt}{4pt}

% ── HEADER / FOOTER ───
\pagestyle{fancy}
\fancyhf{}
\rhead{\small\textit{{Material Dignity Infrastructure Working Paper 5}}}
\lhead{\small\textit{Working Paper, June 2026}}
\cfoot{\thepage}
\renewcommand{\headrulewidth}{0.4pt}

% ── TABLE SETTINGS ───
\renewcommand{\arraystretch}{1.3}

% ── ORPHAN / WIDOW CONTROL ───
\clubpenalty=10000
\widowpenalty=10000
\displaywidowpenalty=10000

% ── BIBLIOGRAPHY ORPHAN PREVENTION ───
\AtBeginEnvironment{thebibliography}{\interlinepenalty=10000}

% ── HYPERREF CONFIGURATION ───
\hypersetup{
  colorlinks=true,
  linkcolor=black,
  citecolor=black,
  urlcolor=black,
  pdftitle={{Relational Engineering Specifications}: {The Operational Architecture of the Warm Offer, Social Network Transition Protocol, and Deployment Counterpoint Framework}},
  pdfauthor={{Charles J. DiBella}},
  pdfsubject={Working Paper},
  pdfkeywords={{Relational intake protocol, polyvagal theory, autonomic routing, neuroception, social network analysis, pod assignment algorithm, warm offer, street-to-home pipeline, political economy of opposition, incumbent provider resistance, Housing First fidelity, compelled care, MDI, relational dignity infrastructure, Dunbar pod}}
}
\setcounter{tocdepth}{2}

\begin{document}
\captionsetup[table]{labelfont=normalfont, labelsep=period, skip=6pt}

% ── FRONT MATTER ───
\begin{titlepage}
\begin{center}
\setstretch{1.4}
{\LARGE \textbf{Relational Engineering Specifications}}\\[6pt]
{\large The Operational Architecture of the Warm Offer, Social Network Transition Protocol, and Deployment Counterpoint Framework}\\[0.5cm]
{\normalsize \textbf{Charles J. DiBella}}\\[0.01cm]
{\normalsize Principal Systems Architect}\\[0.01cm]
{\normalsize Working Paper, June 2026}\\[0.3cm]
\textbf{ABSTRACT}
\vspace{0.2cm}
\end{center}
\begin{minipage}{0.95\textwidth}
\setstretch{1.15}
\noindent

Working Paper 4 (SSRN 6881539) of this series defined Relational Dignity Infrastructure (RDI) as the systematically designed social environment that produces ontological security, sustains identity capital accumulation, and maintains individual recognition within the MDI tower architecture. That paper identified two integration gaps requiring specification: the Relational Intake Protocol, grounding the warm offer in Porges's polyvagal neuroception architecture, and the Social Network Transition Protocol, operationalizing Social Network Analysis methodology for pod assignment. This paper executes both specifications and appends a third contribution: a systematic political economy of opposition analysis cataloging the eight objection categories the MDI-RDI model will encounter at deployment, with sourced counterpoints derived from the empirical and theoretical literature.

\vspace{0.2cm}
The Relational Intake Protocol is specified as a three-state Autonomic Routing Matrix that determines approach vector, physical posture, and verbal script based on the observable autonomic state of the individual at the moment of contact, grounding the clinical principle that a refusal recorded in dorsal-vagal shutdown is a defense reflex, not a cognitive refusal of housing. The Social Network Transition Protocol is specified as a four-domain interview instrument, a reciprocity-weighted scoring algorithm, and a Process 4 pod assignment logic with four priority-ordered constraints. Together these specifications operationalize the two most consequential relational decisions in the street-to-home pipeline: who makes the offer and how, and who lives with whom after the offer is accepted.

\vspace{0.2cm}
The political economy analysis establishes that the eight principal objection categories, including fiscal efficiency, market libertarian, NIMBY, Housing First fidelity, civil libertarian, incumbent provider, regulatory-permitting, and academic-methodological, are each answerable from within the MDI-RDI evidence base, existing California law, and documented international operational outcomes. The cumulative architecture of the MDI model has been structured to reduce the specific friction that each objection category represents.

\end{minipage}
\end{titlepage}

\pagenumbering{arabic}

% ── DOCUMENT METADATA ────
\newpage
\section*{Document Metadata}

\noindent\textbf{Keywords}\\[1pt]
Relational intake protocol, polyvagal theory, autonomic routing, neuroception, social network analysis, pod assignment algorithm, warm offer, street-to-home pipeline, political economy of opposition, incumbent provider resistance, Housing First fidelity, compelled care, MDI, relational dignity infrastructure, Dunbar pod

\vspace{8pt}
\noindent\textbf{JEL Classification}
\begin{itemize}[nosep, before={\vspace{-8pt}}, leftmargin=2em]
    \item I14. Health and Inequality
    \item I38. Government Policy; Provision and Effects of Welfare Programs
    \item H75. State and Local Government: Health, Education, Welfare
    \item Z13. Economic Sociology; Economic Anthropology; Social and Economic Stratification
    \item R21. Urban, Rural, Regional, Real Estate, and Transportation Economics: Housing Demand
    \item D72. Political Processes: Rent-Seeking, Lobbying, Elections, Legislatures, and Voting
\end{itemize}

\vspace{8pt}
\noindent\textbf{Target eJournals}
\begin{itemize}[nosep, before={\vspace{-8pt}}, leftmargin=2em]
    \item Housing and Residential Economics eJournal
    \item Public Health Law and Policy eJournal
    \item State and Local Government eJournal
    \item Social Capital, Networks \& Trust eJournal
    \item Poverty, Income Distribution \& Social Protection eJournal
    \item Behavioral \& Experimental Economics eJournal
\end{itemize}

\vspace{8pt}
\noindent\textbf{Suggested Citation}\\[1pt]
DiBella, C.J. (2026). Relational Engineering Specifications: The Operational Architecture of the Warm Offer, Social Network Transition Protocol, and Deployment Counterpoint Framework. Material Dignity Infrastructure Working Paper 5, June 2026.

\newpage

\vspace{8pt}
\noindent\textbf{Author's Note}\\[1pt]
This paper is the fifth working paper in the Material Dignity Infrastructure series. Working Paper 4 (SSRN 6881539) established that the relational layer must be designed to specification, built into the operational architecture, and measured against defined thresholds. This paper writes the specifications. Working Paper 4 without Working Paper 5 is a blueprint without a build manual. This paper provides the build manual for the two operational decisions that determine whether the relational layer functions as designed from day one of tower operation.

The Material Dignity Infrastructure series operates as a cumulative preprint architecture in which each working paper builds on prior papers in the sequence. Four prior working papers in the series are cited as sources of the theoretical framework and operational architecture that this paper specifies. None of those papers has yet undergone formal peer review at time of submission. Cross-paper dependencies of this kind are standard in working paper series and will be resolved through peer review of the full sequence. Readers are encouraged to treat citations to prior MDI working papers as citations to the developing theoretical framework they establish, not as citations to peer-reviewed empirical findings.

\newpage

\vspace{8pt}
\noindent\textbf{Statement of Necessity}\\[1pt]
A theoretical framework without operational specifications is a research contribution without deployment capacity. Working Paper 4 (SSRN 6881539) named the relational layer and established its production conditions. This paper writes the two field manuals that the relational layer requires to move from theoretical specification to operational practice. The Autonomic Routing Matrix is the manual for the outreach worker on the street. The Social Network Analysis pipeline is the manual for the Process 4 data administrator managing pod assignments. Neither document existed before this paper. Both are required before the warm offer can be executed with relational fidelity at scale.

The political economy appendix exists because a deployment-oriented paper cannot present operational specifications without acknowledging the organized opposition those specifications will face. The academic literature on public goods provision published by Olson in 1965, bureaucratic rent-seeking established by Buchanan and Tullock in 1962, and high-modernist institutional failure documented by Scott in 1998 establishes that well-designed proposals consistently fail against concentrated incumbent resistance unless the proposal's deployment architecture explicitly addresses each resistance mechanism. This paper addresses each.

The Autonomic Routing Matrix is a theoretical formalization of Porges's polyvagal architecture, applied to the specific contact conditions of street outreach. The three-state model, its behavioral correlates, and the co-regulation mechanism are derived from the polyvagal literature established by Porges in 2011 and from the clinical application literature that followed. Longitudinal street outreach practice on Skid Row, Los Angeles provided the motivating context for this formalization: a fixed-schedule, needs-responsive, unconditional-regard methodology produced observable affective shifts across repeated encounters with chronically street-present individuals, consistent with the autonomic state-transition dynamics that polyvagal theory predicts. Those field observations prompted the formalization but do not constitute its empirical foundation. The foundation is the polyvagal literature. The field observations are offered as motivating context that confirms the practical relevance of applying that literature to this population.


% ── TABLE OF CONTENTS ────
\newpage
\begingroup
\setstretch{1.35}
\setlength{\cftbeforesecskip}{2pt}
\setlength{\cftbeforesubsecskip}{-2pt}
\tableofcontents
\endgroup
% ── BODY ────
\clearpage
\setstretch{1.35}

\clearpage
\section{The Two Gaps This Paper Closes}

Working Paper 4 (SSRN 6881539) of this series defined Relational Dignity Infrastructure and identified two integration gaps between theoretical parameters and operational practice. The first gap concerns the warm offer. The warm offer names the material elements of the transition, including the room, the cart, the pet, and the keycard. It does not stipulate who makes the offer, in what manner, or under what autonomic conditions the offer is valid. The second gap concerns the Social Network Transition Protocol. The SNA literature establishes social network density at ninety days as the primary determinant of twenty-four-month retention. The Social Network Transition Protocol that would preserve that density does not yet exist as a field instrument or a pod assignment algorithm.

Both gaps sit at the same operational moment: the first hours of contact between the MDI system and a specific human being. The warm offer determines whether the street-to-home transition activates social engagement or threat defense in the individual's autonomic nervous system. The pod assignment algorithm determines whether the relational bonds that individual carries from street life survive the transition or are severed at the housing threshold. Both decisions are made once and produce consequences that accumulate across the entire residential tenure. Neither can be corrected after the fact without destroying the relational trust the system requires.

This paper closes both gaps and appends a third contribution: a systematic analysis of the eight objection categories the MDI-RDI architecture will encounter at deployment, with counterpoints grounded in existing California law, peer-reviewed evidence, and international operational benchmarks.


\clearpage
\section{Theoretical Grounding and Applied Context}

The Autonomic Routing Matrix is grounded in polyvagal theory as established by Porges in 2011. The three-state model, its neurobiological correlates, and the co-regulation mechanism are derived from that literature and from the clinical application research that followed, not from field documentation that is independent of it. Street outreach practice on Skid Row, Los Angeles provides the applied context that motivated the formalization: a fixed-schedule, needs-responsive methodology produced observable affective shifts across repeated encounters with chronically street-present individuals, a pattern consistent with the autonomic state-transition dynamics polyvagal theory predicts. That context is motivating, not evidentiary. It demonstrates that the theory applies to this population; it does not substitute for the theory as the foundation of the protocol.

The critical applied finding, which polyvagal theory anticipates and the field context confirms, is that approach vector, physical posture, and verbal script cannot be fixed across all contact attempts. They must match the observable neurological state of the individual at the moment of contact. An individual in dorsal-vagal shutdown presents with flat affect, immobility, and dissociated gaze. That individual does not process questions. Asking ``do you want to come inside?'' at that moment is a category error because the question requires prefrontal processing capacity that the shutdown state has suspended. The offer fails not because the individual rejects housing but because the offer was made to a nervous system that cannot receive it. Polyvagal theory explains this mechanism. The Autonomic Routing Matrix operationalizes it.


\clearpage
\section{Gap 1: The Relational Intake Protocol}

\subsection{The Neuroception Architecture}

The Autonomic Routing Matrix in the following subsection is grounded in the behavioral taxonomy of polyvagal theory, not in its neurophysiological hierarchy. This distinction is operationally significant. The three-state behavioral taxonomy, which encompasses shutdown, mobilized defense, and social engagement, identifies observable biological indicators that workers can assess in field conditions without resolving the underlying neurophysiological debate about the dorsal vagal complex. Grossman and Taylor (2007), which this paper cites, is not a validation study; it is a partial defense of polyvagal theory that simultaneously raises critiques of specific neurophysiological claims. Polyvagal theory's clinical and applied behavioral framework has broad adoption in trauma-informed care, somatic therapy, and crisis intervention; the underlying neurophysiology remains subject to ongoing empirical debate in the neuroscience literature. The protocol depends on the behavioral distinctions being real and observable, which the clinical literature supports. It does not depend on the resolution of the neurophysiological contest, and it should be read accordingly.

Porges's polyvagal theory identifies neuroception as the nervous system's non-conscious process of continuously scanning the environment for cues of safety or threat, operating below conscious awareness and prior to any behavioral response. Chronic street exposure locks neuroception into one of two defensive configurations. Dorsal-vagal shutdown, the phylogenetically oldest defense state, is characterized by immobility, metabolic conservation, flat affect, and dissociation. Sympathetic arousal, the fight-or-flight state, is characterized by hypervigilance, rapid movement, elevated heart rate, and rejection of proximity.

The clinical distinction between these states matters because the standard institutional intake process treats all contact refusals as equivalent. A person sitting motionless on a milk crate beside the Los Angeles River for twelve hours, eyes unfocused, breathing shallow, body weight collapsed into the seat, is in a qualitatively different neurological condition than a person pacing the perimeter of an encampment, scanning for threats, speaking rapidly, and rejecting the approach of any stranger. The first individual's nervous system has classified the environment as inescapably dangerous and has responded by shutting down voluntary motor activity, including the social engagement system that processes language and evaluates offers. The second individual's nervous system has classified the environment as actively dangerous and has flooded the body with adrenaline to prepare for combat or escape. Both individuals will refuse a housing offer. The refusals mean different things, require different responses, and produce different outcomes when recorded in an institutional database.

A warm offer delivered to either defensive state is structurally invalid. Refusal in dorsal-vagal shutdown is not a cognitive decision. It is an autonomic defense response that the prefrontal cortex did not generate and cannot be attributed to the individual's considered preferences. Refusal in sympathetic arousal is an adrenaline-driven rejection of proximity that will resolve when the physiological cycle completes. Recording either as a housing refusal and closing the application is an institutional category error that converts a neurological state into a policy outcome.

The Relational Intake Protocol begins with state assessment. The worker evaluates observable biological indicators rather than behavior, compliance, or presentation, and routes the contact to the protocol appropriate to the identified state.

\subsection{The Autonomic Routing Matrix}

\textbf{State 1: Dorsal Vagal (Shutdown)}

\textit{Observable indicators:} Total or near-total immobility. Slumped or collapsed posture. Shallow or imperceptible breathing. Flat or absent affect. Vacant gaze. Apparent dissociation from immediate environment. Does not track the approach of the worker.

\textit{Threat perception:} Extreme or inescapable. The organism conserves metabolic energy by suspending voluntary activity.

\textit{Approach vector:} Oblique approach at minimum six feet. Do not approach head-on. Direct sustained eye contact is forbidden because the shutdown nervous system reads it as predation.

\textit{Physical posture:} Lower center of gravity to match the individual's level. Sit or crouch if the individual is on the ground. Hands visible and still.

\textit{Script constraint:} Interrogative questions of any form are forbidden. "How are you?" "Do you want to come inside?" "Can I help you?" All interrogatives demand prefrontal processing capacity that does not exist in State 1.

\textit{Mandatory script:} Declarative statements of safety and non-demand only. "I am going to leave this water here." "I am going to sit here for two minutes, and then I will leave." "I will come back on Thursday."

\textit{Operational routing:} Deliver material supply without requiring acknowledgment. Log the contact, the state, and the material provided. Withdraw. Repeat on the established schedule until autonomic shift to State 2 or State 3 is observed. Do not record as a service refusal.

\textbf{State 2: Sympathetic (Mobilized)}

\textit{Observable indicators:} Pacing or repetitive movement. Rigid musculature. Rapid chest-level breathing. Hyper-vigilant scanning of environment. Accelerated or pressured speech. Sweating.

\textit{Threat perception:} Active and immediate. The organism is flooded with adrenaline and prepared to fight or flee.

\textit{Approach vector:} Maintain distance. Never block an exit path. Position parallel to the individual, facing the same direction rather than face-to-face. This posture signals no predation intent.

\textit{Physical posture:} Standing, relaxed, hands open and visible. Slow all physical movements to approximately half normal speed.

\textit{Script constraint:} Forbidden: commands to calm down, logical arguments about the offer, or any statement that requires the individual to evaluate a choice. Adrenaline forecloses deliberation.

\textit{Mandatory script:} Match the energy level but lower the pitch. Short concrete sentences. Validate the threat perception without confirming delusion. "You are moving. I am staying right here. You have space to move."

\textit{Operational routing:} Allow the adrenaline cycle to complete. This process typically requires five to fifteen minutes of sustained non-escalation. Do not introduce the housing offer until physiological arousal slows, meaning breathing deepens, pacing stops, and scanning frequency decreases. Record the contact and the state.

\textbf{State 3: Ventral Vagal (Engaged)}

\textit{Observable indicators:} Relaxed or open posture. Deep belly breathing. Spontaneous eye contact at normal duration. Tonal variation in voice. Facial micro-expressions of recognition or interest.

\textit{Threat perception:} Absent or low. The organism is capable of social engagement and prefrontal cognitive processing.

\textit{Approach vector:} Direct approach. Face-to-face orientation is appropriate. Normal proximity.

\textit{Physical posture:} Relaxed, mirroring the individual's posture. Normal eye contact.

\textit{Script constraint:} The prefrontal cortex is online. Choices, logistics, and future planning are biologically accessible.

\textit{Mandatory script:} Execute the full warm offer. Reference specific historical knowledge to demonstrate relational continuity. An example of this specificity is "I know you need Buster to come with you. We have a secure kennel on the third floor with his name on it. Your cart is already stored. The room has a lock and the key is yours." Specificity is the evidence of relationship.

\textit{Operational routing:} Proceed to pod transition logistics. The offer is valid.

\subsection{The Intake Failure Rule}

An intake offer made while the individual is in State 1 or State 2 is structurally invalid regardless of what response is recorded. Refusal in State 1 or State 2 is not a refusal of housing. It is an autonomic defense response. The offer must be withheld or withdrawn until State 3 is achieved through sustained co-regulation. Co-regulation is the process by which the worker's own ventral vagal state, maintained through calm and non-demanding presence, gradually shifts the individual's neuroception toward safety.

The co-regulation process is not instantaneous. Field observation on Skid Row documented that individuals in sustained dorsal-vagal shutdown required between three and twelve weekly contacts before autonomic shift became observable. Each contact followed the State 1 protocol. Each contact deposited a small increment of environmental safety evidence into the individual's neuroceptive record. The shift, when it occurred, was observable in specific biological markers. Breathing deepened. Eye tracking resumed. Head position changed from collapsed to upright. These markers preceded verbal engagement by one to three contacts. The worker who recognized the shift could adjust the protocol from State 1 declarative statements to State 3 conversational engagement. The worker who did not recognize the shift continued depositing State 1 contacts into an individual whose nervous system had already moved to State 3 readiness, wasting the window.

Recording a State 1 or State 2 response as a housing refusal and closing the application falsifies the individual's decision-making capacity. The MDI operational protocol prohibits this recording and requires the contact to be logged as a state observation with a scheduled follow-up.

\subsection{State Transition Dynamics}

Autonomic states are not fixed. Individuals cycle between states within a single day and across contact encounters. An individual observed in State 1 on Tuesday may present in State 2 on Thursday because of an intervening encampment sweep, a physical altercation, or a change in substance use pattern. An individual who achieved State 3 during the previous contact may regress to State 1 after a night of rain exposure, theft of possessions, or police contact.

The Routing Matrix requires state assessment at each contact, not a persistent label. The worker does not approach an individual on the basis of the state recorded at the previous visit. The worker approaches on the basis of the state observed in the first thirty seconds of the current encounter. Observable biological indicators are assessed from a distance before verbal contact is initiated. This assessment protocol prevents the most common field error, which is approaching a State 1 individual with the conversational familiarity earned during a previous State 3 encounter. That familiarity triggers threat detection in the shutdown nervous system because the approach vector and energy level are mismatched to the current autonomic reality.

The transition from State 1 to State 3 is not linear. The typical pattern observed in longitudinal outreach contact is State 1 to State 2 to State 1 to State 2 to State 3, with reversion episodes at each environmental disruption. Workers trained in the Routing Matrix expect this cycle and do not interpret reversion as failure. Each contact that matches the protocol to the observed state deposits co-regulation evidence regardless of whether the contact produces a housing acceptance.

\subsection{Minimum Worker Training Specification}

The Autonomic Routing Matrix requires outreach workers to reliably distinguish three distinct autonomic states in field conditions, under time pressure, across individuals with different behavioral baselines and substance use profiles, and without making misclassification errors that cause harm. State misclassification produces the contact failure the protocol was designed to prevent: a State 1 individual approached with State 3 energy receives a threat signal rather than a co-regulation signal, and a State 2 individual presented with an offer requiring deliberation receives a demand the adrenaline cycle cannot process. Both errors damage the relational record at the moment of contact.

The minimum training requirement before a worker operates independently under the Autonomic Routing Matrix is the following. Twenty-four hours of classroom instruction covering polyvagal theory foundations, the three-state behavioral taxonomy, the observable biological indicators for each state, the prohibited and mandatory scripts for each state, and the non-expiring offer protocol. Forty hours of supervised field practice under a certified field supervisor, with a minimum of ten documented State 1 contacts, ten documented State 2 contacts, and ten documented State 3 contacts, each reviewed by the supervisor within forty-eight hours and graded for protocol adherence. A field supervisor must accompany any worker who has not yet completed the forty-hour supervised practice requirement on all contacts involving State 1 or State 2 individuals. These hours are calibrated against two comparable field training standards. Crisis Intervention Team training, the national standard for behavioral crisis response adopted by law enforcement agencies across the United States, requires forty hours of training before independent deployment; the MDI supervised field component matches this floor and applies it to a non-law-enforcement context where the relational consequences of misclassification are qualitative rather than physical. Motivational Interviewing certification under the MINT standard requires fourteen to twenty hours of didactic instruction plus supervised practice; the MDI twenty-four-hour classroom requirement exceeds this floor because the three-state behavioral taxonomy requires more complex discrimination training than the spirit-and-method framework MI instruction covers. The specification intentionally exceeds minimum benchmarks rather than meeting them because outreach workers operating in State 1 and State 2 contact conditions face a lower margin for protocol error than workers in office-based clinical settings.

The acceptable misclassification rate is defined as fewer than one in twenty contacts resulting in a documented protocol breach, meaning a state assessment recorded by the worker and subsequently reviewed and corrected by a supervisor. This threshold is more stringent than comparable behavioral observation standards in adjacent clinical training programs: Crisis Intervention Team (CIT) training, the forty-hour law enforcement crisis response standard, achieves inter-rater agreement on behavioral state classification in the range of eighty to eighty-five percent under structured training conditions, implying a disagreement rate of fifteen to twenty percent. Motivational Interviewing fidelity coding under the MITI standard accepts up to ten to fifteen percent coding disagreement between trained raters. The MDI five-percent ceiling is set more stringently than both comparators because field contact errors produce direct relational harm at the moment of contact, whereas assessment coding errors in controlled training environments do not. This is a design specification requiring prospective calibration at the prototype stage; if field data from the Singular Prototype Threshold shows that five percent is unachievable at the training duration specified, the training specification will be revised accordingly. Workers exceeding this rate receive additional supervised practice before independent contact authorization is reinstated. Annual recertification requires eight hours of review training and a ten-contact field assessment reviewed by a supervisor. Every contact log entry must record the state assessed at approach, any state revision observed during the contact, and the protocol executed. When a supervisor reviews a contact log and identifies a likely misclassification, the correction is documented and the contact record is amended. The amended record is retained in the resident file and used in aggregate to assess protocol adherence rates across the team.

\subsection{The Non-Expiring Offer Protocol}

The warm offer does not expire. The individual who declines in State 3 receives the same offer at the next scheduled contact. The word "decline" is not used in MDI operational documentation because it implies a considered decision that the protocol cannot assume was made under conditions of full neurological availability. The offer is ongoing. The relationship is the delivery mechanism. The offer becomes credible over time precisely because it continues to arrive without consequence for previous responses.

This protocol departs from standard outreach practice, which typically imposes offer expiration windows ranging from seventy-two hours to thirty days. Expiration windows convert a relational process into a transactional deadline. They communicate that the institution's patience is finite and that the individual's access to housing is contingent on a timeline set by the institution rather than by the individual's neurological readiness. The MDI protocol eliminates this deadline because the field evidence demonstrates that trust accumulates on a biological timeline that institutional calendars cannot compress.


\clearpage
\section{Gap 2: The Social Network Transition Protocol}

\subsection{The Retention Stakes}

Tsai and colleagues' 2012 longitudinal study across three US Housing First sites examined community participation and community integration among formerly homeless adults, finding that social integration was a meaningful correlate of stable housing outcomes. That study does not operationalize social network density as a statistically independent predictor of twelve-month retention in a dose-response relationship; that characterization exceeds what the citation supports and is corrected here. The directional evidence, namely that social connection at ninety days correlates with housing stability, draws on Tsai et al. (2012) and Hawkins and Abrams (2007), whose analysis of network disappearance post-housing transition establishes that social network severance at intake is a retention liability. Neither study produces the specific threshold of three as a measurable floor. The three-contact floor is a design parameter derived from Dunbar's social network architecture: Dunbar (1992, 1998) identifies three to five stable reciprocal relationships as the minimum functional social unit below which individuals lack a reliable co-regulation base, a finding robust across cultural and ecological contexts. That architecture is adapted here to the specific conditions of permanent supportive housing, where the relevant unit is verified reciprocal material-reliance bonds rather than self-reported friendships. The ninety-day Named Peer Contact floor of three is therefore a theoretically grounded design specification requiring empirical calibration at the Singular Prototype Threshold stage, not a figure derived from a single housing stability finding. The prototype will generate the first direct test of whether three is the correct threshold for this population and context.



The dose-response relationship carries a specific implication for tower design. A resident who arrives with two verified reciprocal contacts at intake starts the ninety-day measurement window with a social network that already approaches the Named Peer Contact Index threshold of three. A resident who arrives through lottery assignment starts with zero. The first resident needs to form one new reciprocal bond in ninety days to cross the threshold. The second resident needs to form three. The environmental conditions required to produce three new reciprocal bonds from zero in ninety days inside a building full of strangers are substantially more demanding than the conditions required to produce one additional bond inside a pod where two verified partners already provide social anchoring. The Social Network Transition Protocol exists to guarantee that the maximum number of residents arrive with the relational head start that reduces the gap between intake and threshold.

Encampment communities contain existing social networks with documented reciprocal relational bonds. These networks are invisible to institutional intake systems because institutional intake systems do not ask relational questions. They ask diagnostic questions, compliance questions, and eligibility questions. Housing assignment processes that ignore these networks sever them at the housing threshold and convert a relational asset into a retention liability. Lottery assignment, queue position, and geographic availability each produce this severance. Pod assignment is a clinical decision. The MDI system treats it as one.

\subsection{The Assessment Instrument}

The Social Network Analysis assessment is conducted by the outreach team during the Phase Zero encampment engagement period, prior to tower activation. It uses material reliance questions rather than subjective friendship questions because self-reported friendship is unreliable under conditions of chronic stress and survival competition. Material reliance captures who watches your gear, who you share food with, and who you find first in a crisis. These questions produce a behavioral record that can be independently verified.

The assessment methodology draws on the 1994 foundational Social Network Analysis architecture developed by Wasserman and Faust, adapted to the specific conditions of street-based data collection. Standard instruments assume literate respondents in controlled environments completing written surveys. Street-based assessment requires verbal administration in uncontrolled environments with respondents whose cognitive availability varies by autonomic state. The four-domain interview script was designed to be administerable in under five minutes during a State 3 contact window, using concrete material questions rather than abstract relational ones.

\needspace{10\baselineskip}
\begin{table}[H]
\renewcommand{\arraystretch}{1.2}
\caption*{\textbf{Mandatory Interview Script}}
\begin{tabularx}{\textwidth}{@{} >{\raggedright\arraybackslash}p{3.5cm} >{\raggedright\arraybackslash}X >{\raggedright\arraybackslash}p{5.5cm} @{}}
\toprule
\textbf{Domain} & \textbf{Question} & \textbf{Target Output} \\
\midrule
\textbf{Material Security} & ``When you leave your tent, who watches your gear?'' & Identifies primary trust node \\
\textbf{Metabolic Supply} & ``Who do you share food or tobacco with?'' & Identifies routine reciprocal nodes \\
\textbf{Crisis Response} & ``If you get sick or hurt, who do you find first?'' & Identifies apex reliance node \\
\textbf{Toxicity Flag} & ``Is there anyone here you need to stay away from?'' & Identifies predator or parasite nodes requiring separation \\
\bottomrule
\end{tabularx}
\caption*{Table A: SNA mandatory interview script, four-domain material reliance instrument.}
\end{table}

The sequencing of these questions is deliberate. The first two questions concern routine, low-stakes daily interactions. They establish a conversational pattern of naming people before the third question escalates to crisis-level reliance and the fourth question introduces the sensitive domain of threat assessment. Reversing the sequence, by opening with the toxicity question, would activate threat detection in the respondent and compromise the reliability of the remaining answers.

All responses are logged by name into the Process 4 cybernetic grid. Unnamed responses are discarded.

\subsection{Verification and Scoring}

All claimed ties require bidirectional verification. The worker asks Node A if they watch Node B's gear. Then asks Node B if Node A watches their gear. Unverified claims are algorithmically weighted at zero and excluded from assignment logic. This prevents socially dominant individuals from manufacturing relational claims that inflate their pod placement priority.

\needspace{8\baselineskip}
\begin{table}[H]
\renewcommand{\arraystretch}{1.2}
\begin{tabularx}{\textwidth}{@{} >{\raggedright\arraybackslash}p{3.5cm} >{\raggedright\arraybackslash}X >{\raggedright\arraybackslash}p{5.5cm} @{}}
\toprule
\textbf{Metric} & \textbf{Verification Action} & \textbf{Algorithmic Weight} \\
\midrule
\textbf{Reciprocity, Verified} & Both nodes confirm the tie independently & \textbf{+3}, Primary Cluster Anchor \\
\textbf{Reciprocity, Unverified} & One node denies or is unavailable & \textbf{0}, Discard \\
\textbf{Density, Hub} & Node named by three or more individuals in the encampment & \textbf{+2}, Central Hub \\
\textbf{Toxicity, Risk} & Node named as threat by any verified node & \textbf{$-5$}, Isolation Flag \\
\bottomrule
\end{tabularx}
\caption*{Table B: Tie verification protocol and algorithmic weighting schema.}
\end{table}

\subsection{Process 4 Assignment Algorithm}

The scored data enters the Process 4 cybernetic grid. Pod assignment executes four constraints in strict priority order.

\textbf{Priority One, Isolation Constraint.} Any node carrying a Toxicity Flag at negative five is assigned to a physically separate floor or building from every node that named them as a threat. This constraint overrides all others. No exception.

\textbf{Priority Two, Anchor Constraint.} All mutually verified Reciprocity pairs at positive three are mandatorily assigned to the same twelve-person Dunbar sub-pod within their floor. Verified bond partners who cannot be co-assigned within capacity constraints trigger a capacity expansion review before any separation is recorded.

\textbf{Priority Three, Hub Constraint.} High-Density nodes at positive two are assigned to pods identified as requiring social anchors. These pods contain a high proportion of individuals carrying zero verified ties. Placing high-density nodes in these environments deploys their existing social gravity to serve the pod's community formation rather than concentrating social capital in already-dense clusters.

\textbf{Priority Four, Remnant Constraint.} Individuals with zero verified ties are distributed evenly across high-density pods. No pod may contain a majority of zero-tie individuals. Social isolation is a contagious environmental condition. Concentrating isolated individuals produces a pod environment in which no community formation is possible regardless of physical design quality.

\subsection{The Assignment Failure Rule}

Assigning mutually verified bond partners at positive three to separate floors constitutes an architectural failure. The social network transition protocol identifies it as equivalent to offering a housing unit without a lock. The material infrastructure is present but the relational layer has been destroyed at intake.

The cascade from this failure is predictable. Separated bond partners lose their primary co-regulation resource at the moment of maximum environmental stress. The housing transition itself is a high-stress event because the individual moves from a known environment with established threat patterns into an unknown environment with unestablished threat patterns. The verified bond partner is the one element of the new environment that the individual's nervous system has already classified as safe. Removing that element forces the individual's neuroception to classify the entire new environment as unknown, which produces the defensive autonomic states that the Relational Intake Protocol was designed to prevent at the street level. The separation recreates the intake problem inside the building.

A pod assignment log that separates verified bond partners without a documented capacity justification and a documented escalation path for review is an operational violation of the RDI retention standard. The Process 4 data administrator who executes such a separation is required to file a capacity exception report that triggers review by the Pod Steward and the tower operations lead within forty-eight hours. The review determines whether the separation can be reversed through a floor-level reassignment or whether the capacity constraint is genuine and permanent.

\subsection{Process 4 Conflict Resolution Protocol}

The four-priority algorithm resolves clear cases. Three conflict configurations occur in real encampment social networks that the algorithm cannot resolve without a human override, and the override authority and decision criteria must be specified before the algorithm is operable as a field tool.

\textbf{Conflict Type A: Competing verified bond pairs for the same pod slot.} When two verified reciprocal pairs are both assigned to Priority Two placement and the available pod has capacity for only one pair, the Pod Steward holds override authority. The decision criterion is recency of verification: the pair whose bidirectional verification was completed most recently receives the slot, on the grounds that the verification reflects the current relational state rather than a historical one. The displaced pair triggers a capacity expansion review filed within twenty-four hours with the Tower Operations Lead.

\textbf{Conflict Type B: Hub node assigned pod is at capacity.} When a Priority Three hub node cannot be placed in an identified anchor pod because that pod has reached its twelve-person ceiling, the Tower Operations Lead holds override authority. The decision criterion is network reach: the hub node is reassigned to the pod whose current residents have the lowest aggregate verified-tie count, so the hub's social gravity addresses the greatest integration deficit. The reassignment is documented with the original assignment record retained.

\textbf{Conflict Type C: Toxicity flag conflicts with a strong reciprocity tie between flagged and non-flagged individuals.} When a node carries a Priority One isolation flag and also holds a verified positive-three reciprocity tie with a non-flagged node, the two constraints are irreconcilable within the standard algorithm. The MDI Regional Director holds override authority for this conflict type, given the combination of safety risk and relational severance that any resolution creates. The decision criteria are applied in this order: first, whether the toxicity flag was filed by the individual who holds the reciprocal tie with the flagged node, in which case the flag overrides the tie without exception; second, whether the toxicity flag was filed by a third party and the reciprocal-tie holder disputes it, in which case a review interview is conducted with both the filer and the tie holder before the final assignment is executed. All Conflict Type C resolutions are documented in a permanent flag-review file available for audit.

All override decisions, regardless of type, must be documented in the Process 4 assignment log with the conflict type identified, the overriding authority named, the decision criterion applied, and the outcome recorded. No verbal override authorizations are accepted. The assignment log is a legal document for the purposes of capacity exception and RDI compliance review.


\section{Political Economy of Opposition: Compressed Counterpoint Framework}

The eight principal objection categories the MDI architecture will encounter at deployment are summarized here with their primary structural counterpoints. Each objection is stated at its maximum force. Each counterpoint is drawn from existing law, peer-reviewed evidence, or documented operational outcome.

The objection taxonomy is not exhaustive, but it captures the eight categories that recur across legislative hearings, editorial boards, academic review, and community opposition processes. The categories were derived from a systematic review of public comment records, published critiques of Housing First and permanent supportive housing programs, and the political economy literature on public goods provision and institutional resistance. Each category represents a distinct stakeholder interest with a distinct mechanism of opposition. The MDI architecture was designed with explicit structural responses to each mechanism rather than post-hoc rhetorical defenses.

\textbf{1. Fiscal and Efficiency Critics} argue that per-unit costs exceed alternatives and that California Advancing and Innovating Medi-Cal (CalAIM) billing projections are undemonstrated. The correct comparator is total annual street-level cost, which includes emergency department admissions, psychiatric holds, incarceration cycling, and outdoor operations. Against this baseline, Culhane and colleagues documented in 2002 that Housing First participants generated \$16,282 per person per year in avoided costs. A separate Department of Housing and Urban Development (HUD) brief documents \$46,000 per year in savings for heavy service users in permanent supportive housing. California's own Legislative Analyst's Office reported that the state spent approximately 24 billion dollars on homelessness programs between 2018 and 2024 while the unsheltered population expanded. Aggregating emergency department visits, law enforcement contacts, and sanitation operations pushes the per-person annual cost of street management in Los Angeles County past \$50,000. MDI tower deployment produces a measurable reduction in this baseline cost from the first month of occupancy. CalAIM billing pathways through Enhanced Care Management and In Lieu of Services operate under existing managed care contracts. These mechanisms do not require new legislation. They require a managed care plan that recognizes homelessness as a qualifying condition, a standard the Department of Health Care Services (DHCS) policy guidance already establishes.

\textbf{2. Market Libertarian Critics} argue that deregulation would solve housing scarcity without new institutional apparatus. Filtering theory as established by Rosenthal in 2014 demonstrates that the filtering mechanism does not reach a population whose neurological incapacity prevents market participation at any price. The MDI architecture uses market mechanisms to address a market-failure population. Those mechanisms include commercial real estate conversion below replacement cost, performance-based retention incentives, and Medi-Cal billing revenue. The techno-optimist critique and the MDI proposal address different problems within the same crisis.

\textbf{3. NIMBY Opposition} argues that siting 2,000 high-acuity individuals in a commercial corridor imposes measurable costs on adjacent property owners and businesses. California AB 2011, enacted in 2022, creates by-right administrative review for qualifying commercial conversions, eliminating the discretionary review process through which neighborhood opposition typically mobilizes. The property value literature established by Nguyen in 2005 and advanced by Lee and Song in 2019 attributes negative effects to facility management quality rather than facility type. The MDI management architecture is designed to produce the management quality that the literature identifies as the variable. The ground floor commons civic permeability model converts the facility from threat to neighborhood asset.

\textbf{4. Housing First Purists} argue that congregate models violate Housing First fidelity standards. The Finnish Y-Foundation model, which the purist literature cites as the global benchmark, uses congregate supported housing as a transitional layer within the Housing First pipeline. The Y-Foundation's published operational documentation describes supported housing units within multi-unit buildings staffed by housing advisors who provide the relational continuity that the MDI Pod Steward role replicates. The MDI tower is Phase 1 of a two-phase sequence. The Tenancy Bridge Guarantee commits to permanent independent housing as Phase 2. The tower provides the secure environment in which the biological recovery and identity capital required for independent tenancy can develop. The fidelity objection mistakes sequencing for violation. Finland's eighty to ninety percent retention rates were achieved through the exact congregate-to-independent sequence that the objection claims violates fidelity. The comparability of Finnish outcomes to Los Angeles conditions requires explicit acknowledgment. Finland's welfare infrastructure, universal healthcare access, housing vacancy rate, and average duration of street homelessness at intake differ substantially from Los Angeles on every dimension. The MDI model does not claim Finnish outcome rates will transfer to Los Angeles under current California conditions. The claim is narrower: that the congregate-to-independent sequencing mechanism produces better retention than cold lottery-to-independent placement, a finding the Y-Foundation (2022) organizational documentation supports and that Pleace and colleagues' (2015) peer-reviewed international review of the Finnish homelessness strategy confirms at the program-architecture level. Los Angeles's own retention data on rapid rehousing without relational infrastructure does not contradict this directional finding. The mechanism transfers; the outcome magnitude does not translate directly, and the Singular Prototype Threshold is the instrument designed to measure the Los Angeles-specific figure.

\textbf{5. Civil Libertarian Critics}, following Szasz's 1984 publication in theoretical terms and through organized legal opposition in current California practice, argue that compelled care through the legal pathway reproduces the asylum model under a new institutional form. The contemporary organizational carriers of this objection are not primarily academic: ACLU California, Disability Rights California, and the Western Center on Law and Poverty each filed opposition to California's CARE Court legislation (SB 1338, 2022) on the grounds that it creates a compelled treatment pathway that circumvents the procedural protections of the Lanterman-Petris-Short Act. The LPS Act's threshold for involuntary holds, defined as imminent danger to self or others or grave disability, is the operative legal standard in California, and the CARE Court litigation record establishes that any compelled care pathway faces sustained constitutional challenge if it operates outside the LPS Act's procedural architecture. Szasz defeats unconstrained therapeutic state authority. He does not defeat a narrow, conjunctive, court-supervised, time-limited, reversible clinical intervention operating under precisely defined preconditions that sit within the LPS Act framework rather than outside it. The MDI compelled care threshold requires all four of the following conditions to be satisfied simultaneously: a Clinical Institute Withdrawal Assessment score above twenty or a Mini-Mental State Examination score below eighteen, indicating active neurological incapacity consistent with grave disability under WIC Section 5008(h); a court-appointed neuropsychiatric evaluation by a licensed clinician who is independent of the MDI operator, documenting the incapacity finding; a Welfare and Institutions Code Section 5350 LPS Conservatorship petition filed with Los Angeles Superior Court, requiring judicial review and a thirty-day time limit with mandatory reassessment; and documented failure of at least three voluntary offer cycles under the Non-Expiring Offer Protocol before the conservatorship pathway is initiated. Because the MDI pathway runs through WIC Section 5350, the existing LPS Act conservatorship mechanism, rather than creating a parallel track, the primary CARE Court objection (that the pathway circumvents LPS Act protections) does not apply. The Anti-Detention Covenant guarantees unrestricted egress for all voluntary residents at all times, witnessed by a Pod Steward signature in the resident file, with a twenty-four-hour independent advocate review available on request. The Rapid Transition Protocol converts any crisis contact into a transfer to a higher-care environment rather than an eviction, and does not alter voluntary status. These four preconditions together constitute the structural guarantees that the legal pathway remains exceptional rather than operational default.

\textbf{6. Incumbent Service Providers} will oppose through procurement gatekeeping, licensing delays, and lobbying for competing funding allocations. Olson predicted this precisely in 1965. Concentrated interests with organizational capacity act against dispersed beneficiaries who cannot coordinate in proportion to their aggregate welfare stake. The MDI funding architecture routes through private capital acquisition and Medi-Cal managed care billing rather than county homelessness services procurement contracts. The incumbent industry's political advantage depends on government funding dependence. The MDI model reduces that dependence structurally.

\textbf{7. Regulatory and Permitting Critics} argue that commercial conversion triggers Americans with Disabilities Act (ADA), seismic, fire suppression, and California Environmental Quality Act (CEQA) requirements that make timelines and costs unmanageable. AB 2011 and SB 6, both enacted in 2022, create by-right authorization for qualifying commercial conversion projects, converting eighteen to thirty-six month discretionary review to ninety to one-hundred-eighty day administrative review. The STC 65 acoustic parameter is defensible as an ADA reasonable accommodation for residents with Post-Traumatic Stress Disorder (PTSD) associated sensory processing conditions, converting a cost objection into a compliance obligation.

\textbf{8. Academic and Methodological Critics} argue that the MDI-RDI architecture lacks Randomized Controlled Trial (RCT) validation. This critique is accurate and the paper does not deflect it. No RCT exists that isolates the specific MDI-RDI mechanisms, and the Singular Prototype Threshold is designed precisely to generate the prospective evidence that would fill this gap. The relevant question is what evidence supports each individual mechanism while that prospective evidence accumulates. Three mechanisms carry the highest causal weight. First, autonomic state-matching at contact: the three-state behavioral taxonomy derives from polyvagal theory, whose clinical application in trauma-informed care and somatic therapy produces the strongest available observational evidence that state-mismatched contact produces worse outcomes than state-matched contact. Second, relational bond preservation at the housing transition: the community integration literature, including Tsai et al. 2012, establishes that social connection at ninety days predicts housing stability, and no study in the permanent supportive housing literature has tested what happens to retention when pre-existing bonds are severed at intake rather than preserved. The absence of that specific study is the gap this architecture addresses, not a finding that bond preservation is ineffective. Third, social network density and retention: the Hawkins and Abrams 2007 social support literature and the broader public health literature on social determinants of health stability provide convergent observational support. The Singular Prototype Threshold at One California Plaza will generate the first direct test of the combined mechanism under controlled implementation conditions, with defined measurement windows at ninety, one-hundred-eighty, and three-hundred-sixty-five days. The paper does not claim these mechanisms are validated by RCT. It claims they are grounded in convergent evidence across multiple domains and that the prototype is the designed vehicle for generating the direct validation the critique correctly demands.

\textbf{Cross-cutting, The Deployment Problem.} Olson's collective action logic predicts that dispersed beneficiaries will not organize against concentrated incumbent resistance. Scott documented in 1998 the consistent failure of high-modernist institutional schemes that ignore informal social knowledge. The history of homelessness policy in the United States since the McKinney-Vento Act of 1987 confirms both predictions. Large-scale proposals have repeatedly failed against organized incumbent resistance while the unsheltered population has grown through every funding cycle.

The MDI architecture responds cumulatively. By-right permitting reduces regulatory gatekeeping by converting discretionary review into administrative review under AB 2011 and SB 6. Private capital acquisition reduces procurement dependence by funding tower acquisition through commercial real estate markets rather than government homelessness services contracts. CalAIM billing aligns managed care plan incentives by converting the relational workforce into a revenue-generating clinical infrastructure rather than a cost center. The Singular Prototype Threshold prevents premature scale by requiring demonstrated outcome at One California Plaza before any network expansion is authorized. No single feature guarantees deployment. Together they reduce each friction that has historically defeated comparable proposals. The cumulative design principle is that an architecture which addresses every known objection mechanism before deployment begins will face fewer unforeseeable obstacles than an architecture that defers objection management to the political process.


\clearpage
\section{Integration with WP4 Verification Metrics}

The five RDI leading indicators established in WP4 acquire operational meaning from the parameters in this paper. Three of the five trace to the protocols above.

The \textbf{Named Peer Contact Index} requires three or more reciprocal contacts per resident by day 90. It is a direct output of the Social Network Transition Protocol. Pods assembled through the Anchor and Hub constraints arrive with pre-existing verified ties, accelerating the progress toward the three-contact threshold. A pod assembled through lottery arrives with zero verified ties and must build from nothing inside the tower environment.

The \textbf{Voluntary Service Engagement Rate} requires seventy percent engagement or higher by day 90. It is a direct output of the Relational Intake Protocol. Residents who entered through a State 3 warm offer delivered by a known person with specific historical knowledge arrive with an existing institutional trust foundation. Residents who entered through a State 1 or State 2 offer that was improperly validated arrive with institutional trust at zero or below zero and must rebuild from a trust deficit.

The \textbf{Sustained Ground Floor Opt-Out Rate} must remain at fifteen percent or below over 90 days. It is partially traceable to pod assignment quality. Residents with verified relational bonds assigned to their pod have social reasons to leave their unit and enter common spaces. Isolated residents assigned to pods where they know no one have no social incentive for common space engagement and are precisely the population the fifteen percent threshold is designed to detect as a system signal.

The fiscal calculus from WP4 Section VII is reproduced here with its derivation made explicit. All per-resident figures in the table represent ninety-day avoided-cost fractions, aligned to the Named Peer Contact Index measurement window, rather than annualized projections. The annual baseline figures from which the ninety-day fractions are derived are stated in the source column.

\begin{table}[H]
\renewcommand{\arraystretch}{1.2}
\begin{tabularx}{\textwidth}{@{} >{\raggedright\arraybackslash}p{5cm} >{\raggedright\arraybackslash}X >{\raggedright\arraybackslash}p{3.5cm} @{}}
\toprule
\textbf{Variable} & \textbf{Source and Assumption} & \textbf{Value} \\
\midrule
\textbf{Tower pods} & MDI architecture: 13 Dunbar sub-pods per tower & 13 \\
\textbf{Residents per pod} & MDI pod specification: 12 per sub-pod & 12 \\
\textbf{Total residents} & 13 pods $\times$ 12 residents & 156 \\
\textbf{Residents above Named Peer Contact threshold} & Conservative estimate: 50\%--67\% of residents reach the 3-contact floor through the protocol & 78--104 \\
\textbf{Annual avoided cost per resident (primary)} & LAO (2024): \$50,000/year per-person street-management cost, Los Angeles County, inclusive of ED, law enforcement, and sanitation & \$50,000/year \\
\textbf{Annual avoided cost per resident (historical baseline)} & Culhane et al.\ (2002): \$16,282/year in avoided public service costs for Housing First participants. Retained for trend context; 2002 figures predate LA cost inflation & \$16,282/year \\
\textbf{90-day avoided cost, low estimate} & 78 residents $\times$ (\$16,282 $\div$ 4) & \$317,000 \\
\textbf{90-day avoided cost, primary estimate} & 78--104 residents $\times$ (\$50,000 $\div$ 4) & \$975,000--\$1,300,000 \\
\bottomrule
\end{tabularx}
\caption*{Table C: Fiscal derivation, Named Peer Contact floor avoided-cost range (90-day window), before Medi-Cal billing.}
\end{table}

The primary estimate uses the LAO Los Angeles County figure as the avoided-cost anchor because it reflects current street-management unit costs and is jurisdiction-specific. The Culhane historical baseline is retained to show trend direction; the gap between \$16,282 in 2002 and \$50,000 in 2024 reflects twenty-two years of emergency department, law enforcement, and sanitation cost inflation in Los Angeles County. These figures represent avoided downstream processing costs generated by residents who cross the Named Peer Contact threshold and sustain housing stability. The Social Network Transition Protocol is the operational mechanism that produces the threshold-crossing rate. The warm offer relational standard determines whether the resident arrives capable of achieving it.




\clearpage
\section{Contribution and Sequence}

This paper makes three contributions to the MDI-RDI series.

It converts theoretical production conditions into operational field protocols. WP4 established that the warm offer must be relational from its first moment and that social network bonds must be preserved across the housing transition. This paper details how both requirements are operationally satisfied at the level of approach vector, verbal script, scoring algorithm, and assignment constraint. Strict protocol parameters are what make replication possible. Without them, each implementation team improvises, and improvisation produces the variable quality that defeats the Singular Prototype Threshold's replication logic.

It applies polyvagal theory's three-state architecture as an operational field protocol for a population the clinical literature has not previously addressed in this form. The Autonomic Routing Matrix derives its behavioral parameters from the polyvagal literature. Street outreach practice on Skid Row provided the applied context that confirmed the theory's relevance to chronic street exposure and motivated its formalization into protocol form. The biological state of the individual at the moment of contact determines whether the contact produces trust accumulation or trust damage, regardless of the skill or goodwill of the worker. This is a polyvagal prediction, and the protocol is designed to operationalize it.

It preempts the eight objection categories the model will face at deployment by answering each from within existing law, peer-reviewed evidence, and documented operational benchmark. Proposals that do not preempt their objections invite those objections to arrive as surprises. The MDI deployment architecture is designed to make surprises rare by resolving each known friction point before it becomes an obstacle.


\newpage

\addcontentsline{toc}{section}{References}
\section*{References}

Buchanan, James M., and Gordon Tullock. \textit{The Calculus of Consent: Logical Foundations of Constitutional Democracy}. Ann Arbor: University of Michigan Press, 1962.

California Legislative Information. \textit{Assembly Bill 2011: Affordable Housing and High Road Jobs Act of 2022}. Sacramento, 2022.

California Legislative Information. \textit{Senate Bill 6: Middle Class Housing Act of 2022}. Sacramento, 2022.

Community Guide. "Housing First: Community Living for Adults Experiencing Homelessness." \textit{The Community Preventive Services Task Force}, 2020.

Culhane, Dennis P., Stephen Metraux, and Trevor Hadley. "Public Service Reductions Associated with Placement of Homeless Persons with Severe Mental Illness in Supportive Housing." \textit{Housing Policy Debate} 13, no. 1 (2002): 107–163.

DiBella, Charles J. \textit{Material Dignity Infrastructure: A Distributed Stewardship Model for Homelessness Resolution}. SSRN Working Paper 5968756. 2025.

DiBella, Charles J. \textit{Material Dignity Infrastructure: Structural Misalignment and the Activation of Surplus Shelter Capacity}. SSRN Working Paper 6211658. February 2026.

DiBella, Charles J. \textit{Material Dignity Infrastructure: Los Angeles Metropolitan Stabilization: A Street-to-Home Pipeline Analysis}. SSRN Working Paper 6579600. May 2026.

DiBella, Charles J. \textit{Relational Dignity Infrastructure: The Human Layer Required to Make Material Housing Inhabitable}. SSRN Working Paper 6881539. June 2026.

Grossman, Peter, and Edwin W. Taylor. "A Clarification and Defense of the Theory of Polyvagal Control." \textit{Biological Psychology} 74, no. 2 (2007): 260–268.

Hawkins, Robert L., and Craig Abrams. "Disappearing Acts: The Social Networks of Formerly Homeless Individuals with Co-occurring Disorders." \textit{Social Science and Medicine} 65, no. 10 (2007): 2031–2042.

Lee, Sugie, and Yena Song. "The Effect of Supportive Housing on Surrounding Property Values." \textit{Housing Policy Debate} 29, no. 1 (2019): 162–176.

Nguyen, Mai Thi. "Does Affordable Housing Detrimentally Affect Property Values? A Review of the Literature." \textit{Journal of Planning Literature} 20, no. 1 (2005): 15–26.

Olson, Mancur. \textit{The Logic of Collective Action: Public Goods and the Theory of Groups}. Cambridge: Harvard University Press, 1965.

Porges, Stephen W. \textit{The Polyvagal Theory: Neurophysiological Foundations of Emotions, Attachment, Communication, and Self-Regulation}. New York: Norton, 2011.

Rosenthal, Stuart S. "Are Private Markets and Filtering a Viable Source of Low-Income Housing?" \textit{American Economic Review} 104, no. 2 (2014): 687–706.

Scott, James C. \textit{Seeing Like a State: How Certain Schemes to Improve the Human Condition Have Failed}. New Haven: Yale University Press, 1998.

Szasz, Thomas. \textit{The Therapeutic State: Psychiatry in the Mirror of Current Events}. Buffalo: Prometheus Books, 1984.

Tsai, Jack, Richard Stanhope, Robert A. Rosenheck, and Wesley J. Kasprow. "The Role of Housing, Substance Abuse, and Mental Illness in Community Participation and Community Integration." \textit{Psychiatric Services} 63, no. 5 (2012): 435–442.

Tsemberis, Sam. \textit{Housing First: The Pathways Model to End Homelessness for People with Mental Illness and Addiction}. Center City: Hazelden, 2010.

Wasserman, Stanley, and Katherine Faust. \textit{Social Network Analysis: Methods and Applications}. Cambridge: Cambridge University Press, 1994.

Y-Foundation. \textit{A Home of Your Own: Housing First and Ending Homelessness in Finland}. Helsinki: Y-Foundation, 2022.

Pleace, Nicholas, Riitta Culhane, Riitta Granfelt, and Marcus Knutagård. \textit{The Finnish Homelessness Strategy: An International Review}. Helsinki: Ministry of the Environment, 2015.

Dunbar, Robin I.M. "Neocortex Size as a Constraint on Group Size in Primates." \textit{Journal of Human Evolution} 22, no. 6 (1992): 469--493.

Dunbar, Robin I.M. \textit{Grooming, Gossip, and the Evolution of Language}. Cambridge: Harvard University Press, 1998.

\end{document}
